\subsection{Аппаратная часть и модуль Ledbar}

\subsubsection{Физическая реализация}
Квадрокоптер Геоскан Пионер может быть оснащён модулем светодиодной навигации. Стандартная конфигурация подразумевает использование адресуемой светодиодной ленты или специального модуля расширения.

\begin{infobox}[Устройство пикселя]
Каждый светодиод состоит из трёх кристаллов:
\begin{itemize}
    \item \textbf{Red (R)} — красный канал.
    \item \textbf{Green (G)} — зелёный канал.
    \item \textbf{Blue (B)} — синий канал.
\end{itemize}
Смешивая эти три цвета в разных пропорциях, мы можем получить любой оттенок радуги. Например, включение красного и зелёного даёт жёлтый цвет.
\end{infobox}

\begin{warnbox}[Питание]
Светодиоды потребляют энергию аккумулятора. При написании алгоритмов стоит помнить, что включение всех светодиодов на полную яркость (белый цвет) создаёт максимальную нагрузку на систему питания, хотя для современных дронов это редко является критичным фактором.
\end{warnbox}
